\documentclass[10pt,a4paper]{article}
\usepackage[utf8]{inputenc}
\usepackage[vietnamese]{babel}
\usepackage{amsmath,amssymb,amsfonts}
\usepackage{geometry}
\geometry{a4paper, left=10mm, right=10mm, top=10mm, bottom=10mm}
\usepackage{tikz}
\usetikzlibrary{shapes,arrows.meta,positioning,calc}
\usepackage[most]{tcolorbox}
\usepackage{array,booktabs,multicol}
\usepackage{fontawesome}
\usepackage{helvet}
\renewcommand{\familydefault}{\sfdefault}
\usepackage{parskip}

\definecolor{primary}{RGB}{112, 48, 160}
\definecolor{teal}{RGB}{0, 128, 128}
\definecolor{orange}{RGB}{230, 108, 44}
\definecolor{green}{RGB}{46, 139, 87}
\definecolor{rose}{RGB}{192, 0, 0}

\pagestyle{empty}

\begin{document}

% --- HEADER BANNER ---
\begin{tcolorbox}[
  enhanced,
  colback=violet!8!white,
  colframe=violet!80!black,
  arc=2.5mm, boxrule=1pt,
  left=3mm, right=3mm, top=2mm, bottom=2mm
]
  \begin{tabular*}{\textwidth}{@{\extracolsep{\fill}} lr}
    {\small\bfseries\color{violet!85!black}\faBook\ SỔ TAY KIẾN THỨC TOÁN 2} & 
    {\small\bfseries\color{teal!85!black}\faGraduationCap\ CHƯƠNG TRÌNH GDPT 2018}
  \end{tabular*}
  \vspace{-1mm}
  \begin{center}
    \rule{0.98\textwidth}{0.6pt}
  \end{center}
  \vspace{-2mm}
  \centering
  {\LARGE\bfseries\color{violet!90!black}\faBookmark\ BÀI 54: LUYỆN TẬP CHUNG CÁC SỐ TRONG PHẠM VI 1000}\\[1mm]
  {\itshape\footnotesize\color{black!75} Củng cố đọc viết, cấu tạo số, so sánh và thứ tự các số trong phạm vi 1000}
\end{tcolorbox}

\vspace{-2mm}

\noindent
\begin{minipage}[t]{0.49\textwidth}\vspace{0pt}
\begin{tcolorbox}[
  enhanced,
  colback=violet!4!white,
  colframe=violet!75!black,
  arc=2mm, boxrule=0.8pt,
  left=2mm, right=2mm, top=1.8mm, bottom=1.8mm,
  title={\footnotesize\bfseries\color{white}\faKey\ 1. Hệ thống hóa cấu tạo số đến 1000}
]
{\footnotesize
\begin{itemize}\setlength{\itemsep}{1pt}\setlength{\leftskip}{-2.5mm}
  \item $10$ đơn vị $= 1$ chục; $10$ chục $= 1$ trăm; $10$ trăm $= 1$ nghìn ($1\,000$).
  \item Cấu tạo số $\overline{abc} = a \times 100 + b \times 10 + c$.
  \item \textbf{Số liền trước} của $n$ là $n - 1$:
  \begin{itemize}\setlength{\itemsep}{0.5pt}\setlength{\leftskip}{-2mm}
    \item Số liền trước của $500$ là $499$.
  \end{itemize}
  \item \textbf{Số liền sau} của $n$ là $n + 1$:
  \begin{itemize}\setlength{\itemsep}{0.5pt}\setlength{\leftskip}{-2mm}
    \item Số liền sau của $999$ là $1\,000$.
  \end{itemize}
\end{itemize}
}
\end{tcolorbox}
\end{minipage}
\hfill
\begin{minipage}[t]{0.49\textwidth}\vspace{0pt}
\begin{tcolorbox}[
  enhanced,
  colback=teal!4!white,
  colframe=teal!75!black,
  arc=2mm, boxrule=0.8pt,
  left=2mm, right=2mm, top=1.8mm, bottom=1.8mm,
  title={\footnotesize\bfseries\color{white}\faSliders\ 2. Kĩ năng trọng tâm cần nắm vững}
]
{\footnotesize
\begin{itemize}\setlength{\itemsep}{1pt}\setlength{\leftskip}{-2.5mm}
  \item \textbf{Kĩ năng 1: Đọc, viết số chính xác:}\\
  Chú ý chữ số $0$ ở hàng chục (linh) và số $5$ (lăm/năm).
  \item \textbf{Kĩ năng 2: Viết số thành tổng:}\\
  $670 = 600 + 70$; $803 = 800 + 3$.
  \item \textbf{Kĩ năng 3: Tìm số trên tia số:}\\
  Mỗi vạch cách đều nhau $1, 10$ hoặc $100$ đơn vị.
  \item \textbf{Kĩ năng 4: So sánh và tìm Max / Min} trong một nhóm các số.
\end{itemize}
}
\end{tcolorbox}
\end{minipage}

\vspace{-1.5mm}
\begin{tcolorbox}[
  enhanced,
  colback=orange!4!white,
  colframe=orange!80!black,
  arc=2mm, boxrule=0.8pt,
  left=2mm, right=2mm, top=1.5mm, bottom=1.5mm,
  title={\footnotesize\bfseries\color{white}\faSitemap\ 3. Sơ đồ tư duy liên kết kiến thức số học phạm vi 1000}
]
\centering
\begin{tikzpicture}[x=1.65cm, y=0.8cm]
  \tikzset{arr/.style={-{Stealth[scale=0.8]}}, every node/.append style={align=center}}
  \node[rectangle, rounded corners=2pt, draw=violet!80!black, fill=violet!15, font=\scriptsize\bfseries] (c) at (0,0) {SỐ TRONG\\PHẠM VI 1000};
  \node[rectangle, rounded corners=2pt, draw=teal!80!black, fill=teal!15, font=\tiny\bfseries] (n1) at (2.2,0.9) {ĐỌC - VIẾT\\Theo hàng};
  \node[rectangle, rounded corners=2pt, draw=orange!80!black, fill=orange!15, font=\tiny\bfseries] (n2) at (2.2,0.3) {CẤU TẠO TỔNG\\$a00 + b0 + c$};
  \node[rectangle, rounded corners=2pt, draw=green!70!black, fill=green!15, font=\tiny\bfseries] (n3) at (2.2,-0.3) {SO SÁNH\\Trăm $\to$ Chục $\to$ Đơn vị};
  \node[rectangle, rounded corners=2pt, draw=red!75!black, fill=red!15, font=\tiny\bfseries] (n4) at (2.2,-0.9) {LIỀN TRƯỚC / SAU\\$n-1$ và $n+1$};

  \draw[arr, thick] (c) |- (n1);
  \draw[arr, thick] (c) |- (n2);
  \draw[arr, thick] (c) |- (n3);
  \draw[arr, thick] (c) |- (n4);
\end{tikzpicture}
\end{tcolorbox}

\vspace{-1.5mm}
\noindent
\begin{minipage}[t]{0.49\textwidth}\vspace{0pt}
\begin{tcolorbox}[
  enhanced,
  colback=green!4!white,
  colframe=green!60!black,
  arc=2mm, boxrule=0.8pt,
  left=2mm, right=2mm, top=1.8mm, bottom=1.8mm,
  title={\footnotesize\bfseries\color{white}\faCheckCircle\ 4. Bài tập thực hành tiêu biểu}
]
{\footnotesize
\textbf{1. Tìm số liền trước, số liền sau:}\\
- Số liền trước của $800$ là $\mathbf{799}$.\\
- Số liền sau của $800$ là $\mathbf{801}$.\\
\textbf{2. Điền số thích hợp vào chỗ chấm:}\\
- $300, 400, 500, \mathbf{600}, \mathbf{700}, 800$.\\
- $420, 430, 440, \mathbf{450}, \mathbf{460}, 470$.\\
\textbf{3. Từ ba chữ số $3, 0, 7$}, viết được số lớn nhất có 3 chữ số là $\mathbf{730}$, số bé nhất là $\mathbf{307}$.
}
\end{tcolorbox}
\end{minipage}
\hfill
\begin{minipage}[t]{0.49\textwidth}\vspace{0pt}
\begin{tcolorbox}[
  enhanced,
  colback=rose!4!white,
  colframe=rose!70!black,
  arc=2mm, boxrule=0.8pt,
  left=2mm, right=2mm, top=1.8mm, bottom=1.8mm,
  title={\footnotesize\bfseries\color{white}\faExclamationTriangle\ 5. Lời khuyên khi làm bài thi}
]
{\footnotesize
\begin{itemize}\setlength{\itemsep}{1.5pt}\setlength{\leftskip}{-2.5mm}
  \item \textbf{Cẩn thận khi tìm số liền trước của số tròn trăm:}\\
  Số liền trước của $700$ là $699$ (không phải $690$).
  \item Chữ số đầu tiên của số có 3 chữ số \textbf{luôn luôn khác 0} (ví dụ không thể viết $037$).
  \item Kiểm tra cẩn thận từng hàng khi làm bài toán lập số từ các chữ số cho trước.
\end{itemize}
}
\end{tcolorbox}
\end{minipage}


\end{document}
